\documentclass[10pt]{mypackage}

\usepackage[uprightgreeks,noDcommand]{kpfonts}
\renewcommand{\coloneq}{\coloneqq}
\renewcommand{\eqcolon}{\eqqcolon}
%\usepackage{dsfont}
%\renewcommand{\mathbb}{\mathds}
%\usepackage{homework}
\usepackage{notes}

\usepackage[ backend=bibtex, style = alphabetic, sorting=ynt ]{biblatex}
\addbibresource{../all_references.bib}

\usepackage{parskip}
\frenchspacing

\fancyhf{}
\fancyhead[R]{Avinash Iyer}
\fancyhead[L]{Notions of Equivalence}
\fancyfoot[C]{\thepage}

\setcounter{secnumdepth}{0}

\begin{document}
\RaggedRight
We will discuss different notions of equivalence between groups.
\section{Quasi-Isometry}%
\begin{definition}
  Recall that a map $f\colon X\to Y$ is called a \textit{quasi-isometry} if there exists some constant $C > 0$ such that for all $x,y\in X$, we have that
  \begin{align*}
    \frac{1}{C}d_X\left( x,y \right)-C \leq d_Y\left( f(x),f(y) \right) \leq Cd_X\left( x,y \right) + C
  \end{align*}
  and there is some map $g\colon Y\to X$ and a constant $K$ such that for all $x\in X$ and $y\in Y$, we have $d_X\left( g\circ f(x),x \right) \leq K$ and $d_Y\left( f\circ g(y),y \right)\leq K$. We write $X\sim_{\operatorname{QI}}Y$.
\end{definition}
\begin{remark}
  We call a map $f\colon X\to Y$ such that there exist constants $C,B > 0$ with
  \begin{align*}
    \frac{1}{C}d_X\left( x,y \right)-B \leq d_Y\left( f(x),f(y) \right) \leq Cd_X\left( x,y \right) + B
  \end{align*}
  a $\left( C,B \right)$-quasi-isometric embedding.
\end{remark}
\begin{definition}
  A metric space $\left( X,d \right)$ is said to be \textit{$\left( C,B \right)$-quasi-geodesic} if there are constants $C,B > 0$ such that for any $x,y\in X$, there is a $\left( C,B \right)$-quasi-isometric embedding $\gamma\colon \left[ t,t' \right]\to \R$ such that $\gamma(t) = x$ and $\gamma(t') = y$.
\end{definition}
\begin{example}
  If $X = \left( V,E \right)$ is a connected graph, then the graph metric on $V$ turns $V$ into a $\left( 1,1 \right)$-quasi-geodesic space.

  When $G$ is a finitely generated group, its word metric turns it into a $\left( 1,1 \right)$-quasi-geodesic space.
\end{example}
\subsection{The Milnor--Schwarz Lemma}%
\begin{proposition}[Basic Milnor--Schwarz Lemma]
  Let $G\curvearrowright \left( X,d \right)$ acts on a $\left( C,B \right)$-quasi-geodesic space. Suppose there is a subset $B\subseteq X$ such that the following hold:
  \begin{enumerate}[(i)]
    \item $\operatorname{diam}\left( B \right) < \infty$;
    \item $\bigcup_{g\in G}g.B = X$;
    \item the set $S = \set{g\in G : g.B'\cap B'\neq \emptyset}$ is finite, where
      \begin{align*}
        B' = \set{x\in X: \text{ there exists $y\in B$ such that $d\left( x,y \right)\leq 2b$}}.
      \end{align*}
  \end{enumerate}
  Then,
  \begin{enumerate}[(i)]
    \item $G$ is generated by $S$;
    \item for all $x\in X$, the map $g\mapsto g.x$ is a quasi-isometry with respect to the word metric $d_S$ on $G$.
  \end{enumerate}
\end{proposition}
\nocite{loh_geometric_group_theory,furman2010surveymeasuredgrouptheory}
\printbibliography
\end{document}
